\chapter{EP01｜常数为什么对应 $\delta$？}

\begin{ideabox}[这一集的一句话]
\textbf{信号在时间上铺得越开，频谱就收得越紧；常数把这种收紧推到极限，于是剩下的只有零频率上的一份权重。}
\end{ideabox}

\section{观众的问题链}

这一集不要求观众懂分布，甚至不要求他见过 $\delta$。它只要求一件事：愿意相信“常数只包含零频率”这个直觉，然后亲眼看着这个直觉在积分面前失效。

问题要一层一层递进，每一层都由上一层逼出来：

\begin{center}
\small
\renewcommand{\arraystretch}{1.32}
\begin{tabularx}{\textwidth}{@{}L{0.8cm}Y L{4.4cm}@{}}
\toprule
层 & 问题 & 观众此刻的心理 \\
\midrule
1 & 一个完全不变的信号，包含频率吗？ & 当然，只有零频率。这还要问？\\
2 & 那就把它代进傅里叶积分算一算。 & 应该能算出来吧……\\
3 & 算不出来。振荡得再厉害也没有抵消。 & 那是不是我算错了？\\
4 & 换个办法：先让它衰减，再让衰减消失。 & 这算不算耍赖？\\
5 & 衰减越弱，频谱越窄越高，可是总面积不变。 & 它没有消失，只是挤到一点上去了。\\
\bottomrule
\end{tabularx}
\end{center}

\begin{center}
\small
\renewcommand{\arraystretch}{1.32}
\begin{tabularx}{\textwidth}{@{}L{3.0cm}Y@{}}
\toprule
观众带着什么进来 & 观众带着什么出去 \\
\midrule
“傅里叶变换 = 看看有哪些频率” & “不衰减的信号，普通积分一定失效；但失效的是积分，不是频率的想法。”\\
“常数就是直流，直流就是零频” & “常数的频谱全部压在 $\omega=0$，并且可以用一个叫 $\delta$ 的对象写出等式。”\\
把 $\delta$ 想成一根“很高的线” & “$\delta$ 描述的不是高度，是权重；它的意义在于它在积分里干什么。”\\
\bottomrule
\end{tabularx}
\end{center}

\section{工具：频率探针，以及它为什么要求衰减}

\subsection{为什么用 $e^{-i\omega t}$ 去乘}

傅里叶变换的核 $e^{-i\omega t}$ 是一个“频率探针”。把它乘到信号上：

\begin{equation}
f(t)\cdot e^{-i\omega t}.
\end{equation}

如果 $f$ 本身是频率为 $\omega_0$ 的纯振荡 $e^{i\omega_0 t}$，乘积就是

\begin{equation}
e^{i\omega_0 t}\cdot e^{-i\omega t}=e^{i(\omega_0-\omega)t}.
\end{equation}

\begin{itemize}[leftmargin=1.6em,itemsep=0.3em]
\item 当 $\omega=\omega_0$ 时，乘积等于 $1$：始终同相，积分不断累加，所以结果很大。
\item 当 $\omega\neq\omega_0$ 时，乘积是一个不停止的旋转，正负两半互相抵消，所以结果很小。
\end{itemize}

这就是“频率分析”的全部直觉：\textbf{用探针去和信号比对，对上了就积累，对不上就抵消}。

\begin{ideabox}[但“抵消”不是一个证明]
初学者最容易在这里滑过去的一句话是：“不匹配的时候，正负抵消，所以积分是零。”

这句话在\textbf{有限窗口}里是对的（而且只是近似），在\textbf{无限区间}里没有意义。原因很朴素：$|e^{-i\omega t}|=1$，被积函数从来没有变小过。“抵消”描述的是振荡的相位在转，而不是数值在衰减。要让它成为一个定理，必须另外给条件。

整个 EP01 的第一处张力就在这里，整个系列的第一处张力也在这里。
\end{ideabox}

\subsection{唯一的充分条件：绝对可积}

\begin{prop}[$\Lone$ 上的傅里叶变换]\label{prop:L1}
若 $f\in\Lone$，则对每个 $\omega$，积分 $\int_{\R}f(t)e^{-i\omega t}\dd t$ 绝对收敛，并且
\begin{equation}
|\fh(\omega)|\le \|f\|_1=\int_{\R}|f(t)|\dd t .
\end{equation}
此外 $\fh$ 在 $\R$ 上一致连续，且 $\fh(\omega)\to0$（$|\omega|\to\infty$）。
\end{prop}

证明只需一句：$|f(t)e^{-i\omega t}|=|f(t)|$，与 $\omega$ 无关，所以积分被一个与 $\omega$ 无关的可积函数控制。连续性与“趋于零”分别是控制收敛与 Riemann--Lebesgue 引理的推论。

\begin{notebox}[这两条性质在画面上的意义]
\begin{itemize}[leftmargin=1.4em,itemsep=0.2em]
\item “$|\fh|\le\|f\|_1$”：频谱有界。所以 EP01 后面出现的峰高 $2/\ep$ 越界增长，恰恰说明 $f_\ep$ 的 $\Lone$ 范数 $2/\ep$ 也在爆炸。两者是同一件事。
\item “$\fh\to0$ 于无穷远”：任何 $\Lone$ 信号的频谱在远处必然衰减到零。所以一条“处处不为零、远处还很亮”的频谱，注定不能来自 $\Lone$ 信号。
\end{itemize}
\end{notebox}

\subsection{常数不在 $\Lone$ 里，但这只说明定理用不上}

$\int_{\R}1\,\dd t=\infty$，所以常数 $f\equiv1$ 不属于 $\Lone$，命题 \ref{prop:L1} 用不上。

\begin{warnbox}[一个非常容易被滥用的推理]
“常数不在 $\Lone$，所以它的傅里叶变换不存在。”

这是\textbf{错的}推理。不在 $\Lone$ 只说明这条定理管不到它，不说明积分一定发散。要断言发散，必须\textbf{具体去算}。

这一集的价值，有一半在于让观众看见“具体去算”这四个字：不是引用定理，而是把 $\int_{-T}^{T}e^{-i\omega t}\dd t$ 真的积出来，然后看着它不收敛。

创作者要能分清这两句话：
\begin{itemize}[leftmargin=1.4em,itemsep=0.2em]
\item 对：常数不在 $\Lone$ 里，所以不能直接用绝对收敛的定理。
\item 错：常数不在 $\Lone$ 里，所以它的积分发散。
\end{itemize}
\end{warnbox}

\subsection{顺带一句：直流为什么是零频率}

$1=e^{i\cdot0\cdot t}$。零频率的意思是“相位不转”，也就是“不随时间变化”。所以“常数”和“零频率”在数学上是同一句话的两种说法。

这句话要在开场说清楚，因为后面所有的“集中在原点”都建立在这个等式上。

\section{直接积分：把失败算给人看}

\subsection{对称截断}

形式上把 $f(t)=1$ 代进定义：

\begin{equation}
\int_{\R}e^{-i\omega t}\dd t .
\end{equation}

为了给它一个机会，先做对称截断（这也是最“讲道理”的截断方式）：

\begin{equation}
I_T(\omega)=\int_{-T}^{T}e^{-i\omega t}\dd t
=\begin{cases}
\dfrac{2\sin(\omega T)}{\omega}, & \omega\neq0,\\[9pt]
2T, & \omega=0 .
\end{cases}\label{eq:IT}
\end{equation}

计算只有一行，但要讲清楚两件事：

\begin{itemize}[leftmargin=1.6em,itemsep=0.35em]
\item $\omega\neq0$ 时：$\left[\dfrac{e^{-i\omega t}}{-i\omega}\right]_{-T}^{T}=\dfrac{e^{-i\omega T}-e^{i\omega T}}{-i\omega}=\dfrac{2\sin(\omega T)}{\omega}$。
\item $\omega=0$ 时：被积函数是 $1$，长度是 $2T$，结果是 $2T$。
\item 两式在 $\omega=0$ 处衔接得上：$\lim_{\omega\to0}\frac{2\sin(\omega T)}{\omega}=2T$。所以这不是一个“有洞”的函数，它在 $\omega=0$ 处是连续的。
\end{itemize}

\begin{ideabox}[顺带记住这个函数]
$I_T(\omega)=\dfrac{2\sin(\omega T)}{\omega}$ 不是凭空出现的配角。它正是\textbf{矩形窗} $w_T=\mathbf{1}_{[-T,T]}$ 的傅里叶变换：
\begin{equation}
\widehat{w_T}(\omega)=\int_{-T}^{T}e^{-i\omega t}\dd t=I_T(\omega).
\end{equation}
这一点在 EP01 里只提一句，但它在 EP06 会变成一个漂亮的闭环：$w_T\to1$（分布意义），而 $\widehat{w_T}\to2\pi\delta$。所以第一集看似“失败”的那条曲线，其实已经在暗中指向答案。\textbf{不要在这一集把它讲成纯粹的废品。}
\end{ideabox}

\subsection{固定 $\omega\neq0$：振幅不变，永不收敛}

固定一个非零的 $\omega$，让 $T\to\infty$。此时

\begin{equation}
I_T(\omega)=\frac{2\sin(\omega T)}{\omega},
\end{equation}

而 $\sin(\omega T)$ 在 $-1$ 与 $1$ 之间永不停歇。振幅 $\dfrac{2}{|\omega|}$ 是一个与 $T$ 无关的常数。所以：

\begin{center}
\textbf{它不是“趋于零”，也不是“趋于无穷”，它是根本没有极限。}
\end{center}

想给观众一个无可争辩的证据，就取两列 $T$：

\begin{equation}
T_n=\frac{\pi/2+2\pi n}{|\omega|},
\qquad
S_n=\frac{3\pi/2+2\pi n}{|\omega|},
\qquad n=0,1,2,\dots
\end{equation}

于是

\begin{equation}
I_{T_n}(\omega)=\frac{2}{|\omega|},
\qquad
I_{S_n}(\omega)=-\frac{2}{|\omega|}.
\end{equation}

两列都趋于无穷，但结果分别停在两个不同的数上。极限不存在。

\begin{scriptbox}[动画要点 01-B]
左图：$I_T(\omega_*)$ 随 $T$ 变化的曲线，取定 $\omega_*=1$，标记点沿 $2\sin(T)$ 上下摆动，\textbf{永不收敛}。

右图：同一时刻的 $I_T(0)=2T$，箭头不断上升。

一句话旁白：\textbf{“非零频率处它一直在摆，零频率处它一直往上冲。这两种都不是‘收敛’。”}
\end{scriptbox}

\subsection{零频率处：不是“很大”，而是“没有上界”}

$\omega=0$ 时 $I_T(0)=2T\to\infty$。注意这里的措辞：不是“频谱在零频率处很高”，而是“随着截断变长，这个量无上界地增长”。

这个区别很重要。正因为它是“无上界地增长”，后面我们才需要一个\textbf{权重}的概念，而不是一个\textbf{高度}的概念。

\subsection{结论怎么写才准确}

\begin{warnbox}[三种说法，只有一种是对的]
\begin{itemize}[leftmargin=1.4em,itemsep=0.3em]
\item \textbf{错}：“非零频率处积分等于零。”——积分不存在，谈不上等于零。
\item \textbf{错}：“零频率处是无穷大，所以频谱是无穷大。”——$\infty$ 不是一个数，不能当频谱值。
\item \textbf{对}：“对固定 $\omega\neq0$，截断积分随 $T$ 持续振荡，没有极限；在 $\omega=0$ 处无上界增长。因此原式不能作为通常的反常积分求值。”
\end{itemize}
\end{warnbox}

\begin{notebox}[一句可以放在旁白里的比喻]
“振荡会抵消”就像说“两支方向相反的队伍拔河，所以绳子不动”。问题是：绳子上根本没有裁判喊停。只要两边一直用力，位置就一直在晃。\textbf{振荡是一种运动，不是一种收敛。}
\end{notebox}

\section{正则化：不是绕开，而是逼近}

\subsection{换一个问题的提法}

直接算 $1$ 失败了。但我们真正想要的不是“算出这个积分”，而是“说清楚常数在频率上的样子”。

于是换一个提法：\textbf{造一列能算的信号，每个都越来越像常数，看它们的频谱收敛到哪里。}

这在数学上叫\textbf{正则化}，在讲故事上叫“先给它一个尾巴，再让尾巴慢慢松开”。

\begin{ideabox}[为什么这不是耍赖]
观众（以及很多创作者自己）会有这个疑问：既然算不出来，为什么换个算得出来的东西算，答案就能算数？

关键区别在于：我们\textbf{没有}声称“$\int e^{-i\omega t}\dd t$ 这个积分等于某个数”。我们声称的是另一件事：\textbf{有一列函数 $f_\ep$，它们逐点趋于 $1$，它们的频谱在一种明确的收敛意义下趋于 $2\pi\delta$，所以我们把 $2\pi\delta$ 定义为 $1$ 的频谱。}

这个“明确的收敛意义”正是 EP04 要做的事。第一集只需要把这件事说成“一种新的、待定义的意义”，并诚实地说“严格定义留到后面”。
\end{ideabox}

\subsection{选哪个外壳}

取

\begin{equation}
f_\ep(t)=e^{-\ep|t|},\qquad \ep>0 .\label{eq:feps}
\end{equation}

三个理由，按重要性排列：

\begin{enumerate}[leftmargin=2em,itemsep=0.4em]
\item \textbf{它绝对可积。}$\|f_\ep\|_1=\int_{\R}e^{-\ep|t|}\dd t=\dfrac{2}{\ep}<\infty$。所以命题 \ref{prop:L1} 适用，积分真的收敛，可以放心算。
\item \textbf{它逐点趋于常数。}对每个固定的 $t$，$e^{-\ep|t|}\to1$（$\ep\to0^+$）。
\item \textbf{它是偶函数且非负。}偶性保证频谱是实偶函数（好画、好检查），非负性保证频谱非负（后面“面积集中”的图像才成立）。
\end{enumerate}

\begin{notebox}[为什么不用单边衰减 $e^{-\ep t}H(t)$]
因为那是“从零跳起来再衰减”，它的极限是阶跃 $H$，不是常数 $1$。这正好是 EP03 的主角。这里要的是\textbf{两边都衰减}，极限才是那条水平线。

顺便记住这个对照：EP01 用双边、EP03 用单边。两个外壳，两个不同的极限对象。
\end{notebox}

\begin{warnbox}[$\ep$ 越小越像常数，但永远不等于常数]
动画结束时千万别说“现在它等于 $1$ 了”。$\ep=0.1$ 时，$f_\ep(10)=e^{-1}\approx0.37$，还差得远。它只是在\textbf{越来越长的时间范围内}接近 $1$。这也是为什么“逐点收敛”必须限定“对每个固定的 $t$”：$t$ 越大，需要的 $\ep$ 越小。
\end{warnbox}

\section{唯一需要跟到底的计算}

这是全片唯一一段要求观众跟着走的代数。它必须拆成三到四步显现，不能一屏出完。

\subsection{第一步：为什么必须拆成两半}

$|t|$ 在 $t=0$ 处换了一张脸：左边是 $-t$，右边是 $t$。所以被积函数在左右半轴是不同的表达式，积分必须分开算。

\begin{equation}
\fh_\ep(\omega)=\underbrace{\int_{-\infty}^{0}e^{\ep t}e^{-i\omega t}\dd t}_{\text{左半轴：}|t|=-t}
+\underbrace{\int_{0}^{\infty}e^{-\ep t}e^{-i\omega t}\dd t}_{\text{右半轴：}|t|=t}.
\end{equation}

\begin{scriptbox}[动画要点 01-D]
左半轴用一种颜色，右半轴用另一种颜色；两种颜色分别“飞向”下面两项分式。合并时两种颜色合成一项。

观众只需要记住一句：\textbf{分母里那个 $\ep^{2}+\omega^{2}$，是左右两半加起来的结果。}
\end{scriptbox}

\subsection{第二步：两个半轴各自算}

\begin{align}
\int_{-\infty}^{0}e^{(\ep-i\omega)t}\dd t
&=\left[\frac{e^{(\ep-i\omega)t}}{\ep-i\omega}\right]_{-\infty}^{0}
=\frac{1}{\ep-i\omega},\\
\int_{0}^{\infty}e^{-(\ep+i\omega)t}\dd t
&=\left[\frac{e^{-(\ep+i\omega)t}}{-(\ep+i\omega)}\right]_{0}^{\infty}
=\frac{1}{\ep+i\omega}.
\end{align}

两处的边界项都要交代清楚，这是“为什么 $\ep>0$ 不能省”的地方：

\begin{itemize}[leftmargin=1.6em,itemsep=0.3em]
\item 左端 $t\to-\infty$：$|e^{(\ep-i\omega)t}|=e^{\ep t}\to0$，因为 $\ep>0$。
\item 右端 $t\to+\infty$：$|e^{-(\ep+i\omega)t}|=e^{-\ep t}\to0$，同样因为 $\ep>0$。
\end{itemize}

如果 $\ep=0$，两侧的模都恒等于 $1$，边界项根本不消失——这正是上一节失败的原因。\textbf{$\ep>0$ 不是技术性假设，它是整个计算能成立的原因。}

\subsection{第三步：通分合并}

\begin{equation}
\fh_\ep(\omega)=\frac{1}{\ep-i\omega}+\frac{1}{\ep+i\omega}
=\frac{(\ep+i\omega)+(\ep-i\omega)}{(\ep-i\omega)(\ep+i\omega)}
=\frac{2\ep}{\ep^{2}+\omega^{2}}.
\end{equation}

记

\begin{equation}
G_\ep(\omega):=\frac{2\ep}{\ep^{2}+\omega^{2}} .\label{eq:Geps}
\end{equation}

\subsection{第四步：检查结果是否说得出道理}

\begin{center}
\small
\renewcommand{\arraystretch}{1.3}
\begin{tabularx}{\textwidth}{@{}L{5.0cm}Y@{}}
\toprule
检查项 & 为什么应当如此 \\
\midrule
$G_\ep$ 是实函数 & $f_\ep$ 是实函数，所以 $\overline{\fh(\omega)}=\fh(-\omega)$；实偶信号的频谱是实偶函数 \\
$G_\ep$ 是偶函数 & $f_\ep$ 是偶函数 \\
$G_\ep\ge0$ & $f_\ep\ge0$ 且偶，频谱可以取到非负值（不是定理，但本例成立） \\
$G_\ep(0)=2/\ep$ & 直接代入；$\ep\to0$ 时爆炸，与“零频率处无上界”呼应 \\
\bottomrule
\end{tabularx}
\end{center}

\begin{ideabox}[一个能让画面变漂亮的视角（可选，不必讲）]
若把 $\omega$ 看成复变量，$G_\ep(\omega)=\dfrac{2\ep}{\ep^{2}+\omega^{2}}$ 有两个极点，位置在 $\omega=\pm i\ep$。

$\ep$ 越小，这两个极点越贴近实轴；当 $\ep\to0$，它们撞上实轴的原点——那正是 $\delta$ “出生”的位置。

这个图像对懂复变的观众非常有效，但\textbf{不要放在主视频里}：它会引入一条全新的解释线，和“面积集中”那条线抢戏。留作补充材料。
\end{ideabox}

\section{三个指标：变高、变窄、面积不变}

只算出一个公式是不够的。要让观众“看见”集中，必须给他三个可以量的数。

\begin{enumerate}[leftmargin=2em,itemsep=0.5em]
\item \textbf{峰高}：$G_\ep(0)=\dfrac{2}{\ep}$。$\ep$ 减半，峰高翻倍。
\item \textbf{半高全宽（FWHM）}：解 $\dfrac{2\ep}{\ep^{2}+\omega^{2}}=\dfrac{1}{2}\cdot\dfrac{2}{\ep}$，即 $\ep^{2}+\omega^{2}=2\ep^{2}$，得 $|\omega|=\ep$。所以宽度是 $2\ep$。$\ep$ 减半，宽度减半。
\item \textbf{总面积}：
\begin{equation}
\int_{\R}G_\ep(\omega)\dd\omega
=2\ep\int_{\R}\frac{\dd\omega}{\ep^{2}+\omega^{2}}
=2\ep\cdot\frac{1}{\ep}\left[\arctan\frac{\omega}{\ep}\right]_{-\infty}^{\infty}
=2\pi .
\end{equation}
\end{enumerate}

\begin{center}
\small
\renewcommand{\arraystretch}{1.35}
\begin{tabular}{@{}cccc@{}}
\toprule
$\ep$ & 峰高 $2/\ep$ & 半高全宽 $2\ep$ & 全实轴面积 \\
\midrule
$1$ & $2$ & $2$ & $2\pi$ \\
$0.5$ & $4$ & $1$ & $2\pi$ \\
$0.2$ & $10$ & $0.4$ & $2\pi$ \\
$0.1$ & $20$ & $0.2$ & $2\pi$ \\
$0.02$ & $100$ & $0.04$ & $2\pi$ \\
\bottomrule
\end{tabular}
\end{center}

\begin{figure}[htbp]
\centering
\begin{tikzpicture}
\begin{axis}[width=0.46\textwidth,height=4.3cm,
  domain=-4:4,samples=200,ymin=0,ymax=1.06,
  xlabel={$t$},ylabel={$f_\ep(t)$},
  title={\small 时域：$\ep$ 越小越平},
  legend style={font=\scriptsize,at={(0.03,0.97)},anchor=north west,draw=none},
  ticklabel style={font=\scriptsize},label style={font=\small}]
\addplot[thick,mainblue] {exp(-1*abs(x))};
\addplot[thick,warnred,dashed] {exp(-0.4*abs(x))};
\legend{$\ep=1$,$\ep=0.4$}
\end{axis}
\end{tikzpicture}\hfill
\begin{tikzpicture}
\begin{axis}[width=0.46\textwidth,height=4.3cm,
  domain=-4:4,samples=400,ymin=0,ymax=5.6,
  xlabel={$\omega$},ylabel={$G_\ep(\omega)$},
  title={\small 频域：$\ep$ 越小越尖},
  legend style={font=\scriptsize,at={(0.97,0.97)},anchor=north east,draw=none},
  ticklabel style={font=\scriptsize},label style={font=\small}]
\addplot[thick,mainblue] {2*1/(1+x^2)};
\addplot[thick,warnred,dashed] {2*0.4/(0.16+x^2)};
\legend{$\ep=1$,$\ep=0.4$}
\end{axis}
\end{tikzpicture}
\caption{同一组 $\ep$ 的时频对应。左图变平，右图变高变窄。注意两侧纵轴刻度不同：右图 $\ep=0.4$ 的峰高是 $5$。}
\end{figure}

\subsection{两个尺度是倒数关系}

时域的特征宽度是 $1/\ep$（$e^{-\ep|t|}$ 衰减到 $e^{-1}$ 所需的 $|t|$ 是 $1/\ep$），频域的特征宽度是 $\ep$。两者相乘等于 $1$。

这不是巧合，它来自尺度变换公式：若 $a>0$，

\begin{equation}
f\!\left(\frac{t}{a}\right)
\;\xleftrightarrow{\ \F\ }\;
a\,\fh(a\omega).\label{eq:scaling}
\end{equation}

验证：令 $t=as$，

\begin{equation}
\int_{\R}f\!\left(\frac{t}{a}\right)e^{-i\omega t}\dd t
=a\int_{\R}f(s)e^{-i\omega a s}\dd s
=a\fh(a\omega).
\end{equation}

取 $a=1/\ep$：时域被拉宽 $1/\ep$ 倍，频域就被压窄 $1/\ep$ 倍，同时幅度被压低 $1/\ep$ 倍。再看 $G_\ep$：

\begin{equation}
G_\ep(\omega)=\frac{2\ep}{\ep^{2}+\omega^{2}}
=\frac{1}{\ep}\cdot\frac{2}{1+(\omega/\ep)^{2}}
=\frac{1}{\ep}\,\widehat{g}\!\left(\frac{\omega}{\ep}\right),
\qquad g(t)=e^{-|t|}.
\end{equation}

与 \eqref{eq:scaling} 完全一致。\textbf{这条公式是给创作者的保险：所有“时间铺开、频率集中”的说法，都能用它一句话验证。}

\begin{warnbox}[不要把尺度对偶说成“不确定性原理”]
可以用“时宽 $\times$ 带宽 $=$ 常数”来打比方，但不要说成“测不准原理”，也不要声称“频谱面积守恒”。

真正的 Plancherel 守恒是 $\|f\|_2^2=\frac{1}{2\pi}\|\fh\|_2^2$，而常数连 $\Ltwo$ 都不属于。EP01 里那个“面积恒为 $2\pi$”，是 $G_\ep$ 这一族特定曲线的性质，不是对所有信号都成立的定律。
\end{warnbox}

\section{从“面积集中”到“取样作用”}

\subsection{把面积标准化}

动画里那条面积为 $1$ 的核，和实际频谱差一个 $2\pi$。把这件事说清楚，否则后面每个系数都会错。

\begin{equation}
K_\ep(\omega):=\frac{G_\ep(\omega)}{2\pi}
=\frac{1}{\pi}\frac{\ep}{\ep^{2}+\omega^{2}},
\qquad
\int_{\R}K_\ep(\omega)\dd\omega=1,
\qquad
G_\ep=2\pi K_\ep .\label{eq:Keps}
\end{equation}

记住这个配对：

\begin{center}
\small
\renewcommand{\arraystretch}{1.3}
\begin{tabularx}{\textwidth}{@{}L{5.4cm}L{1.8cm}Y@{}}
\toprule
对象 & 面积 & 用在哪里 \\
\midrule
$K_\ep$（Cauchy 核 / 归一化核） & $1$ & 动画中的“面积为 $1$ 的尖峰”；取样直觉 \\
$G_\ep$（\textbf{本例实际频谱}） & $2\pi$ & 公式里的 $\F[f_\ep]$ \\
\bottomrule
\end{tabularx}
\end{center}

\subsection{集中性：任意小的邻域，最终装下几乎全部质量}

固定任意 $a>0$（注意：$a$ 固定，然后让 $\ep\to0$）：

\begin{equation}
\int_{-a}^{a}K_\ep(\omega)\dd\omega
=\frac{1}{\pi}\left[\arctan\frac{\omega}{\ep}\right]_{-a}^{a}
=\frac{2}{\pi}\arctan\frac{a}{\ep}
\;\longrightarrow\;\frac{2}{\pi}\cdot\frac{\pi}{2}=1 .
\end{equation}

也就是说：\textbf{不管 $a$ 多小，只要 $\ep$ 足够小，原点这个小区间里就装下了几乎全部的权重。}

\begin{ideabox}[这句话要说得非常小心]
正确的说法：“对每个固定的 $a>0$，$\int_{|x|<a}K_\ep\to1$。”

错误的说法：“$K_\ep$ 除原点外处处趋于零，所以整体趋于零。”——前半句是错的（$K_\ep(0)=1/(\pi\ep)\to\infty$），后半句的推理也是错的。

“几乎每一点都趋于零”与“整体趋于零”是两回事。这个区别就是 EP04 的全部内容，这里只需要\textbf{埋下}它，不需要解释它。
\end{ideabox}

\subsection{取样：为什么会“提取原点值”}

取一个光滑、有界的函数 $\varphi$（后面叫它测试函数，这里只说“一条平滑的曲线”）。要算

\begin{equation}
\int_{\R}K_\ep(\omega)\varphi(\omega)\dd\omega .
\end{equation}

直觉分两段：

\begin{itemize}[leftmargin=1.6em,itemsep=0.35em]
\item \textbf{近处}（$|\omega|$ 与 $\ep$ 同量级）：$\varphi$ 变化很小，几乎等于 $\varphi(0)$；而核在近处的质量几乎等于 $1$。于是贡献约为 $1\cdot\varphi(0)$。
\item \textbf{远处}（$|\omega|\gg\ep$）：核的质量已经趋于 $0$，$\varphi$ 有界也翻不起浪。
\end{itemize}

所以

\begin{equation}
\int_{\R}K_\ep(\omega)\varphi(\omega)\dd\omega\;\longrightarrow\;\varphi(0).\label{eq:sample}
\end{equation}

\begin{scriptbox}[动画要点 01-F]
一层阴影从左右两侧向原点收缩，同时一条 $\varphi$ 曲线叠加在上面。阴影覆盖区域内，$\varphi$ 几乎是常数；最终读数显示 $\to\varphi(0)$。

标记要写“面积 $=1$”，但必须注明是\textbf{全实轴}的面积——可见画框里的阴影永远小于 $1$。
\end{scriptbox}

\subsection{用 $\delta$ 的记号把它写下来}

既然“在积分里提取原点值”就是这个对象唯一的作为，那就给它一个名字：

\begin{equation}
\int_{\R}\delta(\omega)\varphi(\omega)\dd\omega=\varphi(0),
\qquad\text{即}\qquad
\ip{\delta}{\varphi}=\varphi(0).\label{eq:deltadef}
\end{equation}

于是 \eqref{eq:sample} 就可以写成

\begin{equation}
K_\ep\longrightarrow\delta
\qquad(\ep\to0^+),
\end{equation}

而实际频谱是 $G_\ep=2\pi K_\ep$，所以

\begin{equation}
G_\ep=\frac{2\ep}{\ep^{2}+\omega^{2}}\longrightarrow2\pi\delta(\omega),
\qquad
\boxed{\ \F[1](\omega)=2\pi\delta(\omega)\ }.\label{eq:mainresult}
\end{equation}

\section{结论的准确含义}

这一节的每一句话，都是写旁白时的“标准答案”。

\subsection{能说的三句话}

\begin{enumerate}[leftmargin=2em,itemsep=0.5em]
\item \textbf{“等号是在广义函数意义下成立的。”}它不是说那个发散积分突然变成收敛积分；它说的是：这一族频谱在“对测试函数的作用”这个意义上收敛到 $2\pi\delta$。
\item \textbf{“常数的频谱集中在零频率，总权重是 $2\pi$。”}$2\pi$ 是一个权重（面积），不是一个高度。
\item \textbf{“为什么是 $2\pi$？因为逆变换里有一个 $1/(2\pi)$。”}\ 两个系数必须互相配套：正变换不带系数，逆变换带 $1/(2\pi)$，所以正变换在零频率上留下 $2\pi$。
\end{enumerate}

\subsection{不能说的三句话}

\begin{warnbox}[EP01 的三条硬禁令]
\begin{enumerate}[leftmargin=1.4em,itemsep=0.35em]
\item \textbf{不能说}“非零频率处原来的积分等于零”。它不存在。
\item \textbf{不能说}“$\delta$ 在原点无穷高”。它没有“取值”。
\item \textbf{不能说}“因为这族函数除了原点都趋于零，所以整体趋于零”。正好相反：整体不趋于零，它趋于一个集中到原点的对象。
\end{enumerate}
\end{warnbox}

\subsection{为什么“面积”不是“能量”}

有一个很容易犯的错：把 $\int G_\ep=2\pi$ 说成“能量守恒”，进而说“所以 $\delta$ 的能量是 $2\pi$”。

正确的图景是 Plancherel：

\begin{equation}
\|f\|_2^2=\frac{1}{2\pi}\|\fh\|_2^2 .
\end{equation}

对 $f_\ep$：左边 $=\int e^{-2\ep|t|}\dd t=\frac{1}{\ep}$，右边 $=\frac{1}{2\pi}\int G_\ep^2\dd\omega=\frac{1}{\ep}$。两边相等，没问题。

但是 $\ep\to0$ 时，$\|f_\ep\|_2^2=\frac{1}{\ep}\to\infty$。也就是说：\textbf{常数 $1$ 不是有限能量信号，它根本不住在 $\Ltwo$ 里。}所以“能量守恒”这条思路从一开始就不适用于它。

\begin{notebox}[要把这句话留到 EP06]
“$1$ 不在 $\Ltwo$，所以不能用 Plancherel；但它在 $\Sv'$，所以可以用对偶。”这两句要等到 EP06 才说得完整。EP01 只需要说：“这条曲线的面积不是信号能量，别把两者混为一谈。”
\end{notebox}

\subsection{一个独立的自检}

做完这一集，可以用另一条路验证 $\F[1]=2\pi\delta$：先算 $\F[\delta]$。

若 $\F[\delta]=1$，那么由反演（在 EP06 证明）应当有 $\F[1]=2\pi\delta$，因为反演说“变换两次基本回到自己，差一个反射和 $2\pi$”。

\begin{equation}
\ip{\F\delta}{\varphi}=\ip{\delta}{\ph}=\ph(0)=\int_{\R}\varphi(t)\dd t=\ip{1}{\varphi}
\quad\Longrightarrow\quad
\F[\delta]=1 .
\end{equation}

这条验证只需要一句 Fubini，非常干净。它同时也是 EP06 里“$\delta$ 与常数”那一段的预演。

\section{动画需要的关键数值}

\begin{center}
\small
\renewcommand{\arraystretch}{1.35}
\begin{tabularx}{\textwidth}{@{}L{1.4cm}L{2.2cm}L{2.2cm}L{2.4cm}Y@{}}
\toprule
$\ep$ & 峰高 $2/\ep$ & 半高宽 $2\ep$ & 面积 & 画面建议 \\
\midrule
$1$ & $2$ & $2$ & $2\pi$ & 纵轴到 $3$ 即可 \\
$0.5$ & $4$ & $1$ & $2\pi$ & 纵轴到 $5$ \\
$0.2$ & $10$ & $0.4$ & $2\pi$ & 纵轴到 $11$，注意采样加密 \\
$0.1$ & $20$ & $0.2$ & $2\pi$ & 纵轴到 $22$；\textbf{默认不是这个范围} \\
$0.02$ & $100$ & $0.04$ & $2\pi$ & 必须缩放或局部放大 \\
\bottomrule
\end{tabularx}
\end{center}

\begin{warnbox}[三个纯技术但会毁掉一集的坑]
\begin{enumerate}[leftmargin=1.4em,itemsep=0.35em]
\item \textbf{峰顶越界。}$\ep=0.1$ 时峰高 $20$。如果沿用 $\ep=1$ 时的纵轴（到 $3$），整个后半段动画的峰会被裁掉，观众看到的是“一条平线加一根竖线”，完全错误。
\item \textbf{采样不足。}$\ep=0.02$ 时半高宽只有 $0.04$。若在 $[-4,4]$ 上均匀取 $200$ 个点，峰顶根本采不到，画出来是个锯齿。必须用对数分布或在原点附近加密的参数曲线。
\item \textbf{面积标签。}可见画框里的阴影只是全实轴的一部分。标签必须写“全实轴面积 $=2\pi$”，否则观众会拿屏幕上的像素面积去核对，然后发现对不上。
\end{enumerate}
\end{warnbox}

\section{讲解设计：这一集要落空三次}

好的科普不是把正确的答案讲一遍，而是让观众的\textbf{预期}被依次打破。EP01 有三次：

\begin{enumerate}[leftmargin=2em,itemsep=0.6em]
\item \textbf{第一次落空}：直觉明明说“只有零频率”，可是积分算不出来。——观众开始怀疑“是不是我理解错了傅里叶变换”。
\item \textbf{第二次落空}：换了衰减外壳，算出很漂亮的峰；$\ep$ 变小，峰变高变窄，“除了原点都趋于零”。——观众以为这说明“频谱趋于零”，结果发现整体并不趋于零，而是集中。
\item \textbf{第三次落空}：以为 $\delta$ 是“一个特别高的尖峰”，结果发现它的意义根本不在高度，而在“它在积分里取走了 $\varphi(0)$”。
\end{enumerate}

\begin{ideabox}[三次落空各自对应一个动画]
\begin{itemize}[leftmargin=1.4em,itemsep=0.2em]
\item 第一次 $\to$ 01-B：截断积分不收敛。
\item 第二次 $\to$ 01-E：双面板，变高变窄，面积标签不变。
\item 第三次 $\to$ 01-F：阴影向原点集中，读数趋于 $\varphi(0)$。
\end{itemize}
如果某一集剪完发现只有“讲清楚了”，没有“落空感”，那多半是缺了其中一个动画。
\end{ideabox}

\section{讲解禁区（EP01）}

\begin{warnbox}[逐条对照，写旁白时逐句检查]
\begin{enumerate}[leftmargin=1.6em,itemsep=0.28em]
\item 不说“$\delta$ 是原点处无穷高的函数”。
\item 不说“非零频率处原积分 $=0$”。
\item 不说“$K_\ep$ 处处趋于零”（原点是例外，而且趋于无穷）。
\item 不说“因为除了原点都趋于零，所以整体趋于零”。
\item 不说“频谱面积守恒是普遍规律”；只说“本例这条曲线的全实轴面积恒为 $2\pi$”。
\item 不说“这就是信号的能量”；不写 $\|\delta\|^2$ 或 $\delta^2$。
\item 不说“$\ep=0$”或“现在它已经是常数了”。
\item 不混用 $K_\ep$（面积 $1$）与 $G_\ep$（面积 $2\pi$）的标注与纵轴。
\item 不在画面上只画一个频率峰来表示实余弦的频谱（实余弦有正负两个频率）。
\item 不说“把这族函数取极限就得到 $\delta$ 这个函数”——$\delta$ 不是函数，这句要等到 EP05 才说得完整。
\item 不声称“频率分量的大小就是峰高”；对 $\delta$ 而言是权重。
\item 如果缩放纵轴，必须让刻度变化可见，不能让屏幕像素面积冒充数学面积。
\end{enumerate}
\end{warnbox}

\section{与 EP02 的接口}

这一集结束在“$\F[1]=2\pi\delta$”。留下的钩子有两个，第二个才是 EP02 真正用的：

\begin{enumerate}[leftmargin=2em,itemsep=0.45em]
\item \textbf{弱的钩子}：“$\delta$ 到底是不是函数？”——这个问题要到 EP05 才正面回答，中间要吊三集。
\item \textbf{强的钩子}：“如果我把常数乘上 $t$，让它开始线性增长，频谱会怎样？”——这个问题下一集就能回答，而且答案会立刻引出新的疑问（$\delta$ 的导数是什么）。
\end{enumerate}

\begin{scriptbox}[最后 30 秒的旁白骨架]
“我们得到了一件很奇怪的东西：一个不是函数的符号，出现在一个等式的右边。

但更奇怪的是，它其实很好用。下一集，我们把这条水平线掰成一条斜线——然后看看会发生什么。”
\end{scriptbox}

\section{本集回顾卡}

\begin{center}
\small
\renewcommand{\arraystretch}{1.4}
\begin{tabularx}{\textwidth}{@{}L{3.6cm}Y@{}}
\toprule
项目 & 内容 \\
\midrule
主结果 & $\F[1](\omega)=2\pi\delta(\omega)$，广义函数意义 \\
核心计算 & $\F\big[e^{-\ep|t|}\big]=\dfrac{2\ep}{\ep^{2}+\omega^{2}}=G_\ep(\omega)$ \\
三个指标 & 峰高 $2/\ep$；半高全宽 $2\ep$；全实轴面积 $2\pi$ \\
归一化核 & $K_\ep=G_\ep/(2\pi)=\frac{1}{\pi}\frac{\ep}{\ep^{2}+\omega^{2}}$，面积 $1$ \\
安全话术 & “频谱集中在零频率，权重为 $2\pi$” \\
危险话术 & “$\delta$ 是原点无穷高”“非零频率处为零”“面积守恒是普遍规律” \\
下一集的钩子 & 把水平线掰成斜线：$\F[t]$ 会是什么？ \\
\bottomrule
\end{tabularx}
\end{center}

\begin{quizbox}[EP01 自测]
\begin{enumerate}[leftmargin=1.6em,itemsep=0.35em]
\item 从零开始，重算 $f_\ep=e^{-\ep|t|}$ 的傅里叶变换，并说明 $\ep>0$ 在哪里被用到。
\item $G_\ep$ 的半高全宽为什么是 $2\ep$？写出求解过程。
\item 证明 $\int_{\R}G_\ep\dd\omega=2\pi$，并解释为什么这个数不随 $\ep$ 变化。
\item 用 \eqref{eq:scaling} 验证 $G_\ep(\omega)=\frac{1}{\ep}\widehat{g}(\omega/\ep)$。
\item 如果某个动画把纵轴固定在 $[0,3]$，$\ep$ 减小到 $0.1$ 时画面会出什么问题？
\item 用一句话说明“面积不变”与“每点趋于零”为什么可以同时成立，而不矛盾。
\end{enumerate}
\end{quizbox}
